Throughout, $\CAR$ on $n$ modes is realised on $\C^{2^{n}}$ by the Jordan-Wigner matrices $v_k = Z \otimes \cdots \otimes Z \otimes \sigma^{-} \otimes 1 \otimes \cdots \otimes 1$ in the occupation basis $|b_1 \cdots b_n\rangle$, $N_k = v_k^{*}v_k$, and the grading is $\alpha = \Ad(\Gamma)$ with $\Gamma = \prod_k (1 - 2N_k)$, as in Definition \claimref{obj.singularity.grading-parity-automorphism}; a state $\tr(\varrho\,\cdot\,)$ is even if and only if $[\varrho, \Gamma] = 0$. The Killing generator $K_{\mathrm{geo}} = \sum_k\omega_kK_k$ differs from $H_0 := \sum_k\omega_kN_k$ by a constant, so the two generate the same flow. Every configuration $b$ is an eigenvector of $H_0$ with energy $E(b) = \sum_k\omega_kb_k$. The single-mode subalgebra generated by $v_k$ is a copy of $\Mat_2(\C)$; by Theorem \claimref{singularity.kms-uniqueness-factor} its unique KMS state for $\Ad(e^{it\omega_kK_k})$ at $\beta \neq 0$ is the Gibbs state, which has $\omega(v_k) = \omega(v_k^{*}) = 0$ and $\omega(v_k^{*}v_k) = e^{-\beta\omega_k}\omega(v_kv_k^{*})$, and conversely these three relations fix the marginal. At $\beta = 2\pi/\kappa$ and $\omega_k = \kappa k$ the factor is $e^{-2\pi k} = u^{k}$. This is the full reading of (R); the literal reading of Premise \claimref{singularity.r-regularity-fixed-locus}, with $a, b$ in the ladder, imposes only the third relation, so a model satisfying the full reading satisfies the literal one, and a model violating the third relation violates both.

\textbf{(i).} Under (Q) the state is the gauge-invariant quasi-free state with diagonal symbol $n_k$; by Fact \claimref{hawking.imported-quasi-free-states-car-algebra} on a finite window it is the product $\bigotimes_k\operatorname{diag}(1 - n_k, n_k)$, and since (Q) makes it KMS for $K_{\mathrm{geo}}$ at $\beta$, Theorem \claimref{singularity.kms-uniqueness-factor} gives $\varrho_Q = e^{-\beta H_0}/Z$, so $n_k = u^{k}/(1 + u^{k})$. Gauge invariance means invariance under $v_k \mapsto e^{i\theta}v_k$ for all $\theta$; at $\theta = \pi$ this map is $\alpha$, so (P) holds. The restriction of $\varrho_Q$ to the mode-$k$ subalgebra has $\omega(v_k) = 0$ (an odd monomial in a diagonal state) and $\omega(v_k^{*}v_k) = n_k = u^{k}(1 - n_k)$, the three relations above, so (R) holds at each mode. Under (G) the state is faithful and its modular group is $\Ad(e^{itcK_{\mathrm{geo}}})$ for a constant $c \neq 0$; a faithful state is KMS for its modular group at $\beta = -1$, so it is KMS for $\Ad(e^{itK_{\mathrm{geo}}})$ at $\beta' = -c \neq 0$, and by Theorem \claimref{singularity.kms-uniqueness-factor} $\varrho = e^{-\beta' K_{\mathrm{geo}}}/Z$. Since $K_{\mathrm{geo}}$ is a polynomial in the $N_k$ it commutes with $\Gamma$, hence so does $\varrho$: (P). The check confirms the (Q) case on the window $k = -1, 0, 1$ (evenness, stationarity, the three relations at each mode, the occupations, Wick factorisation of $\langle N_iN_j\rangle$ and vanishing hopping and pairing two-point data) and the Gibbs case at the nonresonant weights $(1, 2, 4)$ and $(2, 3, 7)$.

\textbf{(ii).} A state is stationary under the Killing flow if and only if $[\varrho, H_0] = 0$, that is, $\varrho$ preserves every eigenspace $V_E = \operatorname{span}\{|b\rangle : E(b) = E\}$. If no $V_E$ contains configurations of both parities, $\Gamma$ is a scalar $\pm 1$ on each $V_E$, so every $\varrho$ commuting with $H_0$ commutes with $\Gamma$ and the state is even. Two configurations $b, b'$ of equal energy and opposite parity give, with $T = \{k : b_k \neq b'_k\}$ and $\epsilon_k = b_k - b'_k$ on $T$, $\sum_{T}\epsilon_k\omega_k = E(b) - E(b') = 0$ with $|T| \equiv \sum_k b_k - \sum_k b'_k$ odd; conversely an odd signed resonance $(T, \epsilon)$ gives such a pair by putting $b_k = 1, b'_k = 0$ where $\epsilon_k = 1$, the opposite where $\epsilon_k = -1$, and $b = b'$ arbitrary off $T$. So stationarity forces evenness exactly when there is no odd signed resonance (for the converse direction see the constructions next). A zero weight is the resonance $T = \{k\}$; $\omega_a + \omega_b = \omega_c$ is $T = \{a, b, c\}$. Construction with Gibbs marginals, for $0 < u < 1$ and the Gibbs state $\varrho_G = e^{-\beta H_0}/Z$, which has full rank: if $|T| \geq 3$, $\varrho = \varrho_G + c(|b\rangle\langle b'| + |b'\rangle\langle b|)$ with $0 < c \leq p_b = p_{b'}$ (equal populations since equal energies) is positive (it differs from $\varrho_G$ only in the $2 \times 2$ block $\left(\begin{smallmatrix} p & c \\ c & p\end{smallmatrix}\right)$), commutes with $H_0$, is odd, and has the marginals of $\varrho_G$, since the partial trace of $|b\rangle\langle b'|$ onto one mode vanishes when $b$ and $b'$ differ at two or more modes. If $T = \{k\}$ ($\omega_k = 0$) and $m \neq k$ is another mode, $\varrho = \varrho_G + c\,X_kZ_m$ with $X_k = v_k + v_k^{*}$, $Z_m = 1 - 2N_m$: $X_kZ_m$ is odd, commutes with $H_0$ (it changes only $b_k$, at zero energy cost), and its partial trace onto any single mode vanishes ($\tr Z = \tr X = 0$ on the mode it is traced over, with the Jordan-Wigner string contributing further factors $Z$); positivity holds for $|c|$ at most the smallest population, as the explicit instance in (iii) shows. The check computes, for each weight set, the dimension of the commutant of $H_0$, $\sum_E (d_E^{+} + d_E^{-})^{2}$, against that of its even part, $\sum_E ((d_E^{+})^{2} + (d_E^{-})^{2})$, with $d_E^{\pm}$ the numbers of even and odd configurations of energy $E$: $(1,2,4)$: $8, 8$; $(1,2,3)$: $10, 8$; $(-1,0,1)$: $24, 12$; $(-2,\ldots,2)$: $176, 88$; $(1,3,5,7)$: $18, 18$ (its resonance $1 + 7 = 3 + 5$ is even); and it confirms that the two differ exactly when an odd signed resonance exists, on all $35$ three-element weight sets from $\{-3, \ldots, 3\}$ and all $15$ four-element sets from $\{1, \ldots, 6\}$ ($39$ of these $50$ are resonant).

\textbf{(iii).} Model A: weights $0, 1$, $\varrho_G = \tfrac12\id \otimes \operatorname{diag}(1, u)/(1 + u)$, $c = u/(4(1 + u))$. The vectors $(|0\rangle \pm |1\rangle) \otimes |m\rangle$, $m = 0, 1$, are eigenvectors with eigenvalues $(2 + u)/(4(1 + u))$, $u/(4(1 + u))$, $(2 - u)/(4(1 + u))$, $3u/(4(1 + u))$, all positive for $0 < u < 1$, so $\varrho_A > 0$; $\tr\varrho_A = 1$; $[\varrho_A, H_0] = 0$ with $H_0 = N_1$; the three relations hold at both modes ($u^{0} = 1$ at the zero mode, where the marginal is tracial); $\tr(\varrho_A X_0Z_1) = u/(1 + u)$; $[\varrho_A, \Gamma] \neq 0$. The dynamics $H = K_{\mathrm{geo}}$ is number-diagonal, so (L$_E$) holds, and $\varrho_A$ is stationary under it. Model B: weights $1, 2, 3$, $b = 110$, $b' = 001$, both of energy $3$ and population $p = u^{3}/Z$, $c = p/2$; the check confirms trace, stationarity, the relations at the three modes, $\tr(\varrho_B v_3^{*}v_2v_1) \neq 0$, oddness, and positivity by exhibiting $8$ independent eigenvectors with eigenvalues positive on $(0, 1)$. Model B shows that the zero mode is not needed: one three-term relation suffices.

\textbf{(iv).} The state $e^{-2\beta H_0}/Z$ on weights $-1, 0, 1$ is diagonal, hence even and stationary, with diagonal symbol and vanishing pairing data (gauge-invariant quasi-free, Fact \claimref{hawking.imported-quasi-free-states-car-algebra}), and its marginals satisfy $\omega(v_k^{*}v_k) = u^{2k}\omega(v_kv_k^{*})$; at $k = \pm 1$ this differs from $u^{k}\omega(v_kv_k^{*})$ since $u \neq 1$, so (R) fails there, in both readings, while the zero mode is tracial at every period. The check confirms the pattern $[\text{fail}, \text{hold}, \text{fail}]$ at the period $p$ and $[\text{hold}, \text{hold}, \text{hold}]$ at $2p$.

\textbf{(v).} For distinct modes $[H_0, v_i^{*}v_j] = (\omega_i - \omega_j)v_i^{*}v_j$, $[H_0, v_iv_j] = -(\omega_i + \omega_j)v_iv_j$, $[H_0, v_i^{*}v_j^{*}] = (\omega_i + \omega_j)v_i^{*}v_j^{*}$, and $[H_0, N_i] = 0$; the quadratic monomials are linearly independent, so the kernel of $\Ad H_0$ on their span is spanned by the monomials of eigenvalue zero. Because $u \neq 1$ the function $E \mapsto u^{E}$ is injective, so $\varrho_Q = u^{H_0}/Z$ and $H_0$ have the same spectral projections and an operator commutes with $\varrho_Q$ if and only if it commutes with $H_0$. With $\omega_k = \kappa k$ distinct, $\omega_i - \omega_j = 0$ only for $i = j$ and $\omega_i + \omega_j = 0$ only for $\{i, j\} = \{k, -k\}$, $k \neq 0$ ($v_0v_0 = 0$): the stationary quadratic generators are the span of the $N_k$ and of $v_kv_{-k}$, $v_{-k}^{*}v_k^{*}$, as stated. In particular $G$ is stationary, even (quadratic), Hermitian, and $[G, N_1] = -v_1v_{-1} + v_{-1}^{*}v_1^{*} \neq 0$, so it is not in the algebra generated by the $N_k$: it is exactly a pairing term of the kind (L$_E$) excludes (compare Proposition \claimref{singularity.l-stabiliser-tangent}, where such terms are the tangent of the vacuum orbit). The quartic $v_2^{*}v_{-1}^{*}v_1v_0$ has eigenvalue $\kappa(2 - 1 - 1 - 0) = 0$; with its adjoint it commutes with $H_0$, with $\varrho_Q$ and with $N_{\mathrm{tot}}$, and not with $N_2$: a number-conserving hop that moves occupation from the modes $0, 1$ to the modes $-1, 2$ without disturbing the (Q) state. The check verifies these commutators as matrices, the zero-eigenvalue lists on the window $-2, \ldots, 2$, and the eigenvalue rule on three pair monomials of that window.

\textbf{(vi).} $Z_0Z_1 = \operatorname{diag}(1, -1, -1, 1)$ has vanishing partial traces onto either mode, so $\varrho = \varrho_G + cZ_0Z_1$ has the Gibbs marginals; it is diagonal, so even and stationary under every element of the algebra generated by the $N_k$; with $c = u^{3}/(2Z)$, $Z = (1 + u)(1 + u^{2})$, every diagonal entry is positive on $(0, 1)$ (the two lowered ones are $u^{2}(2 - u)/(2Z)$ and $u(2 - u^{2})/(2Z)$). Its connected correlation is $\omega(N_0N_1) - \omega(N_0)\omega(N_1) = c \neq 0$, while a quasi-free gauge-invariant state with diagonal symbol has $\omega(N_0N_1) = n_0n_1$ by Wick's theorem (Fact \claimref{hawking.imported-quasi-free-states-car-algebra}). So (Q) fails. This is an explicit instance of the freedom recorded in Corollary \claimref{singularity.chains-state-class-stability}(3).

\textbf{(vii).} Let $A = v_1^{*}v_{-1}^{*}$ and $\Theta = A - A^{*}$, real and anti-Hermitian. On $\operatorname{span}\{|0_{-1}0_1\rangle, |1_{-1}1_1\rangle\}$ (at either occupation of the zero mode) $\Theta$ acts as a rotation generator and it annihilates the singly occupied pairs, so $\Theta^{3} = -\Theta$ and $e^{\theta\Theta} = 1 + \sin\theta\,\Theta + (1 - \cos\theta)\Theta^{2}$; at $(\cos\theta, \sin\theta) = (\tfrac35, \tfrac45)$ this is a rational orthogonal matrix $U_\theta$. The two configurations it mixes have equal energy $-1 + 1 = 0$ and so equal populations in $\varrho_Q$; hence $[U_\theta, \varrho_Q] = 0$. Let $\varrho_c = \varrho_Q + \tfrac{p}{2}(|1_{-1}0_01_1\rangle\langle 0_{-1}0_00_1| + \text{h.c.})$ with $p$ the common population of the two configurations: it is even, and its single-mode marginals are those of $\varrho_Q$ (the coherence connects configurations differing at two modes), so it lies in the class $\mathcal{C}$ of Corollary \claimref{singularity.chains-state-class-stability}. Conjugating the $2 \times 2$ block $\left(\begin{smallmatrix} p & p/2 \\ p/2 & p\end{smallmatrix}\right)$ by a rotation of angle $\theta$ moves $\mp\tfrac{p}{2}\sin 2\theta$ of population between the two configurations; the check computes $U_\theta\varrho_cU_\theta^{*}$ exactly and finds $\omega(N_1)$ shifted by $-\tfrac{6}{25}\,u/(1 + u)^{2} \neq 0$ (with $p = u/(2(1 + u)^{2})$ and $\sin 2\theta = \tfrac{24}{25}$) and the third relation at $k = 1$ violated. (Coherences of the same sign in both zero-mode blocks would cancel in $\omega(N_1)$ through the Jordan-Wigner string of $v_1$; one block suffices.) So the pairing dynamics leaves $\mathcal{C}$, and the stability clause (2) of that corollary uses (L$_E$)'s pairing exclusion, not only its hopping exclusion.

\textbf{Consequence for the block.} By (i), (Q) implies (P) and (R). By (iii), (P) does not follow from $\{$(R), (L$_E$)$\}$ even with stationarity, and since every model of (Q) is even, the models of (iii) violate (Q). By (iv), (R) does not follow from $\{$(P), (L$_E$)$\}$ with stationarity and quasi-freeness. By (v), (L$_E$) does not follow from $\{$(P), (R), (Q)$\}$ with stationarity. By (vi), (Q) does not follow from $\{$(P), (R), (L$_E$)$\}$ with stationarity. Hence no premise of the block follows from the others except through (Q) $\Rightarrow$ (P), (R), and $\{$(P), (R), (L$_E$)$\}$ is independent. The models are finite and every relation used is verified in exact rational-function arithmetic in $u$; positivity on $0 < u < 1$ (where $u = e^{-2\pi}$ lies) is certified by eigenvectors whose eigenvalues have numerator and denominator polynomials with nonnegative Bernstein coefficients on $[0, 1]$, not all zero.
